% Options for packages loaded elsewhere
% Options for packages loaded elsewhere
\PassOptionsToPackage{unicode}{hyperref}
\PassOptionsToPackage{hyphens}{url}
%
\documentclass[
  11pt,
  ignorenonframetext,
  aspectratio=169,
]{beamer}
\newif\ifbibliography
\usepackage{pgfpages}
\setbeamertemplate{caption}[numbered]
\setbeamertemplate{caption label separator}{: }
\setbeamercolor{caption name}{fg=normal text.fg}
\beamertemplatenavigationsymbolsempty
% remove section numbering
\setbeamertemplate{part page}{
  \centering
  \begin{beamercolorbox}[sep=16pt,center]{part title}
    \usebeamerfont{part title}\insertpart\par
  \end{beamercolorbox}
}
\setbeamertemplate{section page}{
  \centering
  \begin{beamercolorbox}[sep=12pt,center]{section title}
    \usebeamerfont{section title}\insertsection\par
  \end{beamercolorbox}
}
\setbeamertemplate{subsection page}{
  \centering
  \begin{beamercolorbox}[sep=8pt,center]{subsection title}
    \usebeamerfont{subsection title}\insertsubsection\par
  \end{beamercolorbox}
}
% Prevent slide breaks in the middle of a paragraph
\widowpenalties 1 10000
\raggedbottom
\AtBeginPart{
  \frame{\partpage}
}
\AtBeginSection{
  \ifbibliography
  \else
    \frame{\sectionpage}
  \fi
}
\AtBeginSubsection{
  \frame{\subsectionpage}
}
\usepackage{iftex}
\ifPDFTeX
  \usepackage[T1]{fontenc}
  \usepackage[utf8]{inputenc}
  \usepackage{textcomp} % provide euro and other symbols
\else % if luatex or xetex
  \usepackage{unicode-math} % this also loads fontspec
  \defaultfontfeatures{Scale=MatchLowercase}
  \defaultfontfeatures[\rmfamily]{Ligatures=TeX,Scale=1}
\fi
\usepackage{lmodern}

\ifPDFTeX\else
  % xetex/luatex font selection
\fi
% Use upquote if available, for straight quotes in verbatim environments
\IfFileExists{upquote.sty}{\usepackage{upquote}}{}
\IfFileExists{microtype.sty}{% use microtype if available
  \usepackage[]{microtype}
  \UseMicrotypeSet[protrusion]{basicmath} % disable protrusion for tt fonts
}{}
\makeatletter
\@ifundefined{KOMAClassName}{% if non-KOMA class
  \IfFileExists{parskip.sty}{%
    \usepackage{parskip}
  }{% else
    \setlength{\parindent}{0pt}
    \setlength{\parskip}{6pt plus 2pt minus 1pt}}
}{% if KOMA class
  \KOMAoptions{parskip=half}}
\makeatother

\usepackage{color}
\usepackage{fancyvrb}
\newcommand{\VerbBar}{|}
\newcommand{\VERB}{\Verb[commandchars=\\\{\}]}
\DefineVerbatimEnvironment{Highlighting}{Verbatim}{commandchars=\\\{\}}
% Add ',fontsize=\small' for more characters per line
\usepackage{framed}
\definecolor{shadecolor}{RGB}{241,243,245}
\newenvironment{Shaded}{\begin{snugshade}}{\end{snugshade}}
\newcommand{\AlertTok}[1]{\textcolor[rgb]{0.68,0.00,0.00}{#1}}
\newcommand{\AnnotationTok}[1]{\textcolor[rgb]{0.37,0.37,0.37}{#1}}
\newcommand{\AttributeTok}[1]{\textcolor[rgb]{0.40,0.45,0.13}{#1}}
\newcommand{\BaseNTok}[1]{\textcolor[rgb]{0.68,0.00,0.00}{#1}}
\newcommand{\BuiltInTok}[1]{\textcolor[rgb]{0.00,0.23,0.31}{#1}}
\newcommand{\CharTok}[1]{\textcolor[rgb]{0.13,0.47,0.30}{#1}}
\newcommand{\CommentTok}[1]{\textcolor[rgb]{0.37,0.37,0.37}{#1}}
\newcommand{\CommentVarTok}[1]{\textcolor[rgb]{0.37,0.37,0.37}{\textit{#1}}}
\newcommand{\ConstantTok}[1]{\textcolor[rgb]{0.56,0.35,0.01}{#1}}
\newcommand{\ControlFlowTok}[1]{\textcolor[rgb]{0.00,0.23,0.31}{\textbf{#1}}}
\newcommand{\DataTypeTok}[1]{\textcolor[rgb]{0.68,0.00,0.00}{#1}}
\newcommand{\DecValTok}[1]{\textcolor[rgb]{0.68,0.00,0.00}{#1}}
\newcommand{\DocumentationTok}[1]{\textcolor[rgb]{0.37,0.37,0.37}{\textit{#1}}}
\newcommand{\ErrorTok}[1]{\textcolor[rgb]{0.68,0.00,0.00}{#1}}
\newcommand{\ExtensionTok}[1]{\textcolor[rgb]{0.00,0.23,0.31}{#1}}
\newcommand{\FloatTok}[1]{\textcolor[rgb]{0.68,0.00,0.00}{#1}}
\newcommand{\FunctionTok}[1]{\textcolor[rgb]{0.28,0.35,0.67}{#1}}
\newcommand{\ImportTok}[1]{\textcolor[rgb]{0.00,0.46,0.62}{#1}}
\newcommand{\InformationTok}[1]{\textcolor[rgb]{0.37,0.37,0.37}{#1}}
\newcommand{\KeywordTok}[1]{\textcolor[rgb]{0.00,0.23,0.31}{\textbf{#1}}}
\newcommand{\NormalTok}[1]{\textcolor[rgb]{0.00,0.23,0.31}{#1}}
\newcommand{\OperatorTok}[1]{\textcolor[rgb]{0.37,0.37,0.37}{#1}}
\newcommand{\OtherTok}[1]{\textcolor[rgb]{0.00,0.23,0.31}{#1}}
\newcommand{\PreprocessorTok}[1]{\textcolor[rgb]{0.68,0.00,0.00}{#1}}
\newcommand{\RegionMarkerTok}[1]{\textcolor[rgb]{0.00,0.23,0.31}{#1}}
\newcommand{\SpecialCharTok}[1]{\textcolor[rgb]{0.37,0.37,0.37}{#1}}
\newcommand{\SpecialStringTok}[1]{\textcolor[rgb]{0.13,0.47,0.30}{#1}}
\newcommand{\StringTok}[1]{\textcolor[rgb]{0.13,0.47,0.30}{#1}}
\newcommand{\VariableTok}[1]{\textcolor[rgb]{0.07,0.07,0.07}{#1}}
\newcommand{\VerbatimStringTok}[1]{\textcolor[rgb]{0.13,0.47,0.30}{#1}}
\newcommand{\WarningTok}[1]{\textcolor[rgb]{0.37,0.37,0.37}{\textit{#1}}}

\usepackage{longtable,booktabs,array}
\newcounter{none} % for unnumbered tables
\usepackage{calc} % for calculating minipage widths
\usepackage{caption}
% Make caption package work with longtable
\makeatletter
\def\fnum@table{\tablename~\thetable}
\makeatother
\usepackage{graphicx}
\makeatletter
\newsavebox\pandoc@box
\newcommand*\pandocbounded[1]{% scales image to fit in text height/width
  \sbox\pandoc@box{#1}%
  \Gscale@div\@tempa{\textheight}{\dimexpr\ht\pandoc@box+\dp\pandoc@box\relax}%
  \Gscale@div\@tempb{\linewidth}{\wd\pandoc@box}%
  \ifdim\@tempb\p@<\@tempa\p@\let\@tempa\@tempb\fi% select the smaller of both
  \ifdim\@tempa\p@<\p@\scalebox{\@tempa}{\usebox\pandoc@box}%
  \else\usebox{\pandoc@box}%
  \fi%
}
% Set default figure placement to htbp
\def\fps@figure{htbp}
\makeatother





\setlength{\emergencystretch}{3em} % prevent overfull lines

\providecommand{\tightlist}{%
  \setlength{\itemsep}{0pt}\setlength{\parskip}{0pt}}



 


\setbeamertemplate{navigation symbols}{}
\setbeamertemplate{footline}[frame number]
\setbeamertemplate{itemize items}[circle]
\setbeamertemplate{itemize subitem}[triangle]
\setbeamercolor{frametitle}{fg=black}
\setbeamerfont{frametitle}{series=\bfseries}
\newcommand{\indep}{\perp\!\!\!\perp}
\usepackage{booktabs}
% shrink code input (fancyvrb) and code output (verbatim)
\usepackage{fancyvrb}
\fvset{fontsize=\scriptsize}
\makeatletter
\def\verbatim@font{\scriptsize\ttfamily}
\makeatother
\makeatletter
\@ifpackageloaded{caption}{}{\usepackage{caption}}
\AtBeginDocument{%
\ifdefined\contentsname
  \renewcommand*\contentsname{Table of contents}
\else
  \newcommand\contentsname{Table of contents}
\fi
\ifdefined\listfigurename
  \renewcommand*\listfigurename{List of Figures}
\else
  \newcommand\listfigurename{List of Figures}
\fi
\ifdefined\listtablename
  \renewcommand*\listtablename{List of Tables}
\else
  \newcommand\listtablename{List of Tables}
\fi
\ifdefined\figurename
  \renewcommand*\figurename{Figure}
\else
  \newcommand\figurename{Figure}
\fi
\ifdefined\tablename
  \renewcommand*\tablename{Table}
\else
  \newcommand\tablename{Table}
\fi
}
\@ifpackageloaded{float}{}{\usepackage{float}}
\floatstyle{ruled}
\@ifundefined{c@chapter}{\newfloat{codelisting}{h}{lop}}{\newfloat{codelisting}{h}{lop}[chapter]}
\floatname{codelisting}{Listing}
\newcommand*\listoflistings{\listof{codelisting}{List of Listings}}
\makeatother
\makeatletter
\makeatother
\makeatletter
\@ifpackageloaded{caption}{}{\usepackage{caption}}
\@ifpackageloaded{subcaption}{}{\usepackage{subcaption}}
\makeatother

\usepackage{bookmark}
\IfFileExists{xurl.sty}{\usepackage{xurl}}{} % add URL line breaks if available
\urlstyle{same}
\hypersetup{
  pdftitle={Causality and Potential Outcomes},
  pdfauthor={Professor Ji-Woong Chung},
  hidelinks,
  pdfcreator={LaTeX via pandoc}}


\title{\texorpdfstring{Causality and Potential
Outcomes}{Causality and Potential Outcomes}}
\subtitle{\texorpdfstring{BUSS975: Causal Inference in Financial
Research}{BUSS975: Causal Inference in Financial Research}}
\author{Professor Ji-Woong Chung}
\date{}
\institute{Korea University}

\begin{document}
\frame{\titlepage}


\begin{frame}{Outline}
\protect\phantomsection\label{outline}
\begin{enumerate}
\tightlist
\item
  What causality is --- and isn't
\item
  The potential outcomes framework
\item
  Selection bias: the Perfect Regulator
\item
  Randomization and independence
\item
  SUTVA
\item
  Randomization inference
\end{enumerate}
\end{frame}

\section{1. What Causality Is --- and
Isn't}\label{what-causality-is-and-isnt}

\begin{frame}{Motivation}
\protect\phantomsection\label{motivation}
\begin{itemize}
\tightlist
\item
  As researchers, we want to make \textbf{causal} statements:

  \begin{itemize}
  \tightlist
  \item
    What is the effect of a corporate tax change on firms' leverage?
  \item
    What is the effect of CEO equity ownership on risk-taking?
  \item
    What is the effect of a disclosure mandate on the cost of capital?
  \end{itemize}
\item
  We are rarely satisfied with saying variables are merely
  ``associated.''
\item
  But nearly all of our data is \textbf{observational}: nobody
  randomized leverage, ownership, or regulation for us.
\item
  This course is about when --- and how --- observational data can
  answer causal questions.
\end{itemize}
\end{frame}

\begin{frame}{Three goals of empirical analysis}
\protect\phantomsection\label{three-goals-of-empirical-analysis}
\begin{enumerate}
\tightlist
\item
  \textbf{Description} --- how things are or were.

  \begin{itemize}
  \tightlist
  \item
    Has industry concentration increased over time?
  \end{itemize}
\item
  \textbf{Prediction} --- guessing an unknown value, without
  intervening.

  \begin{itemize}
  \tightlist
  \item
    What will this firm's earnings be next quarter?
  \end{itemize}
\item
  \textbf{Causality} --- how \emph{changing} one variable would change
  another, all else equal.

  \begin{itemize}
  \tightlist
  \item
    How would a disclosure mandate change firms' cost of capital?
  \end{itemize}
\end{enumerate}

\vspace{0.5em}

\begin{itemize}
\tightlist
\item
  All three are legitimate. But they require different tools, and
  confusing them is the original sin of empirical work.
\end{itemize}
\end{frame}

\begin{frame}{The common thread: the assignment mechanism}
\protect\phantomsection\label{the-common-thread-the-assignment-mechanism}
\begin{itemize}
\tightlist
\item
  What breaks the link between correlation and causation is \textbf{how
  treatment got assigned}:

  \begin{itemize}
  \tightlist
  \item
    Who chose to be treated? Why? Based on what?
  \end{itemize}
\item
  The central message of this lecture (and this course):
\end{itemize}

\begin{center}
\emph{Correlations are properties of the data.\\
Causal effects are properties of the data \textbf{plus} the assignment mechanism.}
\end{center}

\begin{itemize}
\tightlist
\item
  To formalize this, we need a language for ``what would have happened
  otherwise.''
\end{itemize}
\end{frame}

\begin{frame}{Defining causality: counterfactuals}
\protect\phantomsection\label{defining-causality-counterfactuals}
\begin{itemize}
\tightlist
\item
  Our (narrow) definition: \emph{ceteris paribus}, how does \(D\) affect
  \(Y\)?
\item
  To know the effect of \(D\) on unit \(i\), we need to compare:

  \begin{itemize}
  \tightlist
  \item
    \(Y_i\) \textbf{with} the treatment, and
  \item
    \(Y_i\) \textbf{without} the treatment,
  \item
    with everything else the same.
  \end{itemize}
\item
  One of these is always hypothetical --- a \textbf{counterfactual}.
\item
  \textbf{The Fundamental Problem of Causal Inference} (Holland 1986):
  we can never observe both states for the same unit at the same time.
\end{itemize}
\end{frame}

\section{2. The Potential Outcomes
Framework}\label{the-potential-outcomes-framework}

\begin{frame}{States of the world}
\protect\phantomsection\label{states-of-the-world}
\begin{itemize}
\tightlist
\item
  Imagine each unit exists in two potential states of the world,
  identical in every respect except one: treatment.
\item
  Binary treatment indicator: \[
  D_i = \begin{cases} 1 & \text{if unit } i \text{ is treated} \\ 0 & \text{otherwise} \end{cases}
  \]
\item
  Two \textbf{potential outcomes} for every unit:

  \begin{itemize}
  \tightlist
  \item
    \(Y_i^1\): the outcome unit \(i\) \emph{would} have if treated,
  \item
    \(Y_i^0\): the outcome unit \(i\) \emph{would} have if untreated.
  \end{itemize}
\item
  Both are defined for every unit --- treated or not. That is the trick.
\end{itemize}
\end{frame}

\begin{frame}{The switching equation}
\protect\phantomsection\label{the-switching-equation}
\begin{itemize}
\tightlist
\item
  What we \emph{observe} is only one of the two: \[
  Y_i = D_i\, Y_i^1 + (1 - D_i)\, Y_i^0
  \]
\item
  Treatment ``switches on'' which potential outcome is realized.
\item
  Notice what this equation does \textbf{not} contain: no error term, no
  functional form, no linearity. It is a definition, not a model.
\end{itemize}
\end{frame}

\begin{frame}{Individual treatment effects}
\protect\phantomsection\label{individual-treatment-effects}
\begin{itemize}
\tightlist
\item
  The causal effect of the treatment on unit \(i\): \[
  \delta_i = Y_i^1 - Y_i^0
  \]
\item
  \(\delta_i\) carries an \(i\) subscript.
\item
  Effects are allowed to \textbf{differ across units} from the start.
\item
  The Fundamental Problem, restated: \(\delta_i\) is never observed for
  any unit, because one of its two ingredients is always missing.

  \begin{itemize}
  \tightlist
  \item
    Causal inference is a \textbf{missing data problem} --- but the
    missing value never existed to begin with.
  \end{itemize}
\end{itemize}
\end{frame}

\begin{frame}{A running example: capital injections}
\protect\phantomsection\label{a-running-example-capital-injections}
\begin{itemize}
\tightlist
\item
  Setting: a banking crisis. The government can inject capital into
  banks (think TARP, 2008).
\item
  Outcome \(Y\): bank lending growth (\%) over the following year.
\item
  \(Y_i^1\): bank \(i\)'s lending growth \emph{with} an injection.
\item
  \(Y_i^0\): bank \(i\)'s lending growth \emph{without} one.
\item
  Question: what is the causal effect of capital injections on lending?
\end{itemize}
\end{frame}

\begin{frame}{Ten banks, two potential outcomes}
\protect\phantomsection\label{ten-banks-two-potential-outcomes}
{\def\LTcaptype{none} % do not increment counter
\begin{longtable}[]{@{}cccc@{}}
\toprule\noalign{}
Bank & \(Y^1\) & \(Y^0\) & \(\delta\) \\
\midrule\noalign{}
\endhead
1 & 9 & 2 & 7 \\
2 & 6 & 1 & 5 \\
3 & 7 & 3 & 4 \\
4 & 5 & 2 & 3 \\
5 & 6 & 4 & 2 \\
6 & 7 & 6 & 1 \\
7 & 7 & 7 & 0 \\
8 & 7 & 8 & \(-1\) \\
9 & 7 & 9 & \(-2\) \\
10 & 9 & 8 & 1 \\
\bottomrule\noalign{}
\end{longtable}
}

\begin{itemize}
\tightlist
\item
  Weak banks (low \(Y^0\)) gain a lot; healthy banks gain little --- or
  are even hurt.
\end{itemize}
\end{frame}

\begin{frame}{Average treatment effects}
\protect\phantomsection\label{average-treatment-effects}
\begin{itemize}
\tightlist
\item
  \textbf{Average Treatment Effect (ATE)}: \[
  ATE = E[\delta_i] = E[Y_i^1] - E[Y_i^0]
  \]
\item
  In our table: \(ATE = \frac{7+5+4+3+2+1+0-1-2+1}{10} = 2.0\)
\item
  On average, injections raise lending growth by 2 percentage points.
\item
  But note the heterogeneity: effects range from \(+7\) to \(-2\).

  \begin{itemize}
  \tightlist
  \item
    An average can hide winners and losers.
  \end{itemize}
\end{itemize}
\end{frame}

\begin{frame}{ATT and ATU: different questions, different answers}
\protect\phantomsection\label{att-and-atu-different-questions-different-answers}
\begin{itemize}
\tightlist
\item
  \textbf{Average Treatment effect on the Treated}: \[
  ATT = E[\delta_i \mid D_i = 1]
  \]
\item
  \textbf{Average Treatment effect on the Untreated}: \[
  ATU = E[\delta_i \mid D_i = 0]
  \]
\item
  These are \textbf{different policy questions}:

  \begin{itemize}
  \tightlist
  \item
    ATT: ``Did TARP help the banks that received it?''
  \item
    ATE: ``Would injections help a randomly chosen bank?''
  \item
    ATU: ``Would extending the program to the remaining banks help
    them?''
  \end{itemize}
\item
  When effects are heterogeneous and assignment is not random,
  \(ATT \neq ATE \neq ATU\) --- and each estimator you will meet in this
  course targets a \emph{particular} one of these.
\end{itemize}
\end{frame}

\begin{frame}{What we actually observe}
\protect\phantomsection\label{what-we-actually-observe}
{\def\LTcaptype{none} % do not increment counter
\begin{longtable}[]{@{}ccccc@{}}
\toprule\noalign{}
Bank & \(D\) & \(Y^1\) & \(Y^0\) & \(\delta\) \\
\midrule\noalign{}
\endhead
1 & 1 & 9 & ? & ? \\
2 & 1 & 6 & ? & ? \\
3 & 1 & 7 & ? & ? \\
4 & 1 & 5 & ? & ? \\
5 & 1 & 6 & ? & ? \\
6 & 0 & ? & 6 & ? \\
7 & 0 & ? & 7 & ? \\
8 & 0 & ? & 8 & ? \\
9 & 0 & ? & 9 & ? \\
10 & 0 & ? & 8 & ? \\
\bottomrule\noalign{}
\end{longtable}
}

\begin{itemize}
\tightlist
\item
  Half of the table is missing --- forever. Everything we do from here
  on is a strategy for filling in the question marks \emph{on average}.
\end{itemize}
\end{frame}

\section{3. Selection Bias: The Perfect
Regulator}\label{selection-bias-the-perfect-regulator}

\begin{frame}{The naive comparison}
\protect\phantomsection\label{the-naive-comparison}
\begin{itemize}
\tightlist
\item
  With observational data, the natural move is to compare treated and
  untreated: \[
  SDO = E[Y_i^1 \mid D_i = 1] - E[Y_i^0 \mid D_i = 0]
  \]
\item
  \textbf{Simple Difference in Outcomes} (SDO): mean outcome of the
  treated minus mean outcome of the untreated.
\item
  Every cross-sectional regression of \(Y\) on \(D\) is, at heart, an
  SDO.
\item
  Question: when does SDO recover the ATE?
\end{itemize}
\end{frame}

\begin{frame}{The Perfect Regulator}
\protect\phantomsection\label{the-perfect-regulator}
\begin{itemize}
\tightlist
\item
  Suppose the regulator has perfect foresight: it knows every bank's
  \(\delta_i\) and injects capital into the five banks that \textbf{gain
  the most} (banks 1--5).
\item
  This is not a villain --- it is an \emph{optimizing} regulator,
  allocating scarce capital where it helps most.

  \begin{itemize}
  \tightlist
  \item
    Exactly what a good policymaker (or a good CEO, or a good doctor)
    should do.
  \end{itemize}
\item
  Treated: banks 1--5. Untreated: banks 6--10.
\item
  Now compute the SDO from the observable data.
\end{itemize}
\end{frame}

\begin{frame}{The Perfect Regulator's data}
\protect\phantomsection\label{the-perfect-regulators-data}
\begin{itemize}
\tightlist
\item
  Mean outcome among treated:
  \(E[Y^1 \mid D=1] = \frac{9+6+7+5+6}{5} = 6.6\)
\item
  Mean outcome among untreated:
  \(E[Y^0 \mid D=0] = \frac{6+7+8+9+8}{5} = 7.6\) \[
  SDO = 6.6 - 7.6 = -1.0
  \]
\item
  The truth: \(ATE = +2.0\), and \(ATT = \frac{7+5+4+3+2}{5} = +4.2\).
\item
  The naive comparison gets the \textbf{sign wrong}: injections look
  \emph{harmful} precisely because they were allocated \emph{well}.
\item
  This is the TARP debate in one table: ``bailed-out banks lent less
  than healthy banks'' is an SDO, not a causal effect.
\end{itemize}
\end{frame}

\begin{frame}{The decomposition}
\protect\phantomsection\label{the-decomposition}
\begin{itemize}
\item
  The SDO can always be decomposed --- this is an identity, not an
  assumption (\(\pi\) = share of units treated): \[
  \begin{aligned}
  \underbrace{E[Y^1 \mid D{=}1] - E[Y^0 \mid D{=}0]}_{\text{SDO}}
  = ATE &+ \underbrace{E[Y^0 \mid D{=}1] - E[Y^0 \mid D{=}0]}_{\text{selection bias}} \\
  &+ \underbrace{(1-\pi)\,(ATT - ATU)}_{\text{heterogeneous TE bias}}
  \end{aligned}
  \]
\item
  \textbf{Selection bias}: treated and untreated differ in their
  \emph{untreated} outcomes.
\item
  \textbf{Heterogeneous treatment effect bias}: the treated gain
  differently than the untreated would.
\item
  Memorize this equation. Every research design in this course is an
  attempt to kill these two terms.
\end{itemize}
\end{frame}

\begin{frame}{Proving the decomposition}
\protect\phantomsection\label{proving-the-decomposition}
Let \(\pi = P(D{=}1)\).

\textbf{Step 1.} Add and subtract \(E[Y^0 \mid D{=}1]\) --- the
\emph{counterfactual} mean for the treated: \[
\begin{aligned}
SDO &= E[Y^1 \mid D{=}1] - E[Y^0 \mid D{=}0] \\
    &= \underbrace{E[Y^1 \mid D{=}1] - E[Y^0 \mid D{=}1]}_{ATT} \;+\;
       \underbrace{E[Y^0 \mid D{=}1] - E[Y^0 \mid D{=}0]}_{\text{selection bias}}
\end{aligned}
\]

\textbf{Step 2.} By the law of total expectation,
\(ATE = \pi\,ATT + (1-\pi)\,ATU\). Subtract \(ATE\) from \(ATT\): \[
ATT - ATE \;=\; (1-\pi)\,ATT - (1-\pi)\,ATU \;=\; (1-\pi)(ATT - ATU)
\]

Substituting \(ATT = ATE + (1-\pi)(ATT-ATU)\) into Step 1 gives the
three-term decomposition. \(\blacksquare\)
\end{frame}

\begin{frame}{The decomposition, in our example}
\protect\phantomsection\label{the-decomposition-in-our-example}
\begin{itemize}
\tightlist
\item
  Selection bias: \small \[
  E[Y^0 \mid D{=}1] - E[Y^0 \mid D{=}0] = \frac{2+1+3+2+4}{5} - \frac{6+7+8+9+8}{5} = 2.4 - 7.6 = -5.2
  \] \normalsize

  \begin{itemize}
  \tightlist
  \item
    Treated banks were \emph{much weaker} to begin with.
  \end{itemize}
\item
  Heterogeneous TE bias, with \(\pi = 0.5\): \[
  (1 - 0.5)(ATT - ATU) = 0.5 \times (4.2 - (-0.2)) = +2.2
  \]
\item
  Check: \(\;2.0 - 5.2 + 2.2 = -1.0 = SDO\) \checkmark
\item
  Both bias terms are large; they partially offset, and the net result
  is a sign flip.
\end{itemize}
\end{frame}

\begin{frame}{Selection on gains is endemic in finance}
\protect\phantomsection\label{selection-on-gains-is-endemic-in-finance}
\begin{itemize}
\tightlist
\item
  The Perfect Doctor (Cunningham's version) treats patients who benefit
  most. Finance is full of Perfect Doctors:

  \begin{itemize}
  \tightlist
  \item
    Regulators bail out the banks where intervention matters most.
  \item
    Boards grant equity to the CEOs they expect to respond.
  \item
    Firms go public when going public helps them most.
  \item
    Banks lend to the borrowers they screen as creditworthy.
  \end{itemize}
\item
  Agents \textbf{optimize}. Optimization means treatment is assigned on
  expected gains --- the worst case for naive comparisons.
\item
  This is why finance cannot simply borrow ``big data'' logic: more
  observations do not fix a biased assignment mechanism.
\end{itemize}
\end{frame}

\begin{frame}[fragile]{Simulation: triage vs.~coin flip}
\protect\phantomsection\label{simulation-triage-vs.-coin-flip}
\begin{Shaded}
\begin{Highlighting}[]
\FunctionTok{set.seed}\NormalTok{(}\DecValTok{975}\NormalTok{)}
\NormalTok{n    }\OtherTok{\textless{}{-}} \DecValTok{100000}
\NormalTok{y0   }\OtherTok{\textless{}{-}} \FunctionTok{rnorm}\NormalTok{(n, }\DecValTok{7}\NormalTok{, }\DecValTok{2}\NormalTok{)     }\CommentTok{\# lending growth, no injection}
\NormalTok{gain }\OtherTok{\textless{}{-}} \FunctionTok{rnorm}\NormalTok{(n, }\DecValTok{2}\NormalTok{, }\DecValTok{1}\NormalTok{)     }\CommentTok{\# true individual gain (so ATE = 2)}
\NormalTok{y1   }\OtherTok{\textless{}{-}}\NormalTok{ y0 }\SpecialCharTok{+}\NormalTok{ gain}

\NormalTok{D1 }\OtherTok{\textless{}{-}} \FunctionTok{as.integer}\NormalTok{(y0 }\SpecialCharTok{\textless{}} \FunctionTok{quantile}\NormalTok{(y0, }\FloatTok{0.30}\NormalTok{))  }\CommentTok{\# triage: weakest 30\%}
\NormalTok{D2 }\OtherTok{\textless{}{-}} \FunctionTok{rbinom}\NormalTok{(n, }\DecValTok{1}\NormalTok{, }\FloatTok{0.30}\NormalTok{)                   }\CommentTok{\# coin flip: random 30\%}

\NormalTok{sdo }\OtherTok{\textless{}{-}} \ControlFlowTok{function}\NormalTok{(D) }\FunctionTok{mean}\NormalTok{(y1[D }\SpecialCharTok{==} \DecValTok{1}\NormalTok{]) }\SpecialCharTok{{-}} \FunctionTok{mean}\NormalTok{(y0[D }\SpecialCharTok{==} \DecValTok{0}\NormalTok{])}
\FunctionTok{round}\NormalTok{(}\FunctionTok{c}\NormalTok{(}\AttributeTok{ATE =} \FunctionTok{mean}\NormalTok{(gain), }\AttributeTok{triage =} \FunctionTok{sdo}\NormalTok{(D1), }\AttributeTok{coin =} \FunctionTok{sdo}\NormalTok{(D2)), }\DecValTok{2}\NormalTok{)}
\end{Highlighting}
\end{Shaded}

\begin{verbatim}
   ATE triage   coin 
  2.00  -1.31   1.98 
\end{verbatim}
\end{frame}

\begin{frame}{What the simulation tells us}
\protect\phantomsection\label{what-the-simulation-tells-us}
\begin{itemize}
\tightlist
\item
  The true effect is \(+2\) in both worlds --- same banks, same gains,
  only the \textbf{assignment rule} differs.
\item
  Triage (treat the weakest): SDO is strongly \emph{negative}. An
  observer concludes bailouts destroy lending.
\item
  Coin flip: SDO \(\approx\) ATE. Same data-generating process, correct
  answer.
\item
  With \(n = 100{,}000\): the bias does not shrink with sample size.
  \textbf{Selection bias is not a small-sample problem.}
\end{itemize}
\end{frame}

\section{4. Randomization and
Independence}\label{randomization-and-independence}

\begin{frame}{The independence assumption}
\protect\phantomsection\label{the-independence-assumption}
\begin{itemize}
\tightlist
\item
  What made the coin flip work? Treatment assignment was
  \textbf{independent of potential outcomes}: \[
  (Y^1, Y^0) \indep D
  \]
\item
  In words: knowing a bank's treatment status tells you \emph{nothing}
  about its potential outcomes --- neither its baseline \(Y^0\) nor its
  gain \(\delta\).
\item
  Under independence:

  \begin{itemize}
  \tightlist
  \item
    \(E[Y^0 \mid D{=}1] = E[Y^0 \mid D{=}0]\) \(\Rightarrow\)
    \textbf{selection bias \(= 0\)}
  \item
    \(E[\delta \mid D{=}1] = E[\delta \mid D{=}0]\) \(\Rightarrow\)
    \(ATT = ATU = ATE\) \(\Rightarrow\) \textbf{heterogeneity bias
    \(= 0\)}
  \end{itemize}
\item
  Both terms die, and \(SDO = ATE\). Randomization is the assignment
  mechanism that \emph{guarantees} independence.
\end{itemize}
\end{frame}

\begin{frame}{What randomization balances}
\protect\phantomsection\label{what-randomization-balances}
\begin{itemize}
\tightlist
\item
  Physical randomization balances treated and control groups on:

  \begin{itemize}
  \tightlist
  \item
    observable characteristics (size, leverage, ratings, \ldots),
  \item
    \textbf{unobservable} characteristics (management quality, risk
    appetite, \ldots),
  \item
    potential outcomes themselves, and
  \item
    treatment gains \(\delta_i\).
  \end{itemize}
\item
  No control variables needed. No model needed. The balance is created
  by the \emph{mechanism}, not by the econometrician.
\item
  Practice: in any experiment (or natural experiment), the first table
  is a \textbf{balance table} --- compare pre-treatment covariates
  across groups. Balance on observables is the testable \emph{symptom}
  of a sound mechanism.
\end{itemize}
\end{frame}

\begin{frame}{Where finance experiments live}
\protect\phantomsection\label{where-finance-experiments-live}
\begin{itemize}
\tightlist
\item
  Finance does have real randomized experiments. Four representative
  ones:

  \begin{itemize}
  \tightlist
  \item
    \textbf{Blake, Nosko and Tadelis (2015)} switched eBay's paid search
    off in random markets. Brand-keyword ads had \textbf{no measurable
    benefit} --- traffic substituted to the free organic link --- and
    average returns were negative. Attribution had been measuring
    selection, not causation.
  \item
    \textbf{Karlan and Zinman (2009)} mailed 58,000 South African loan
    offers, randomizing the \emph{offer} rate, the \emph{contract} rate
    and a repayment incentive. The offer rate identifies \textbf{adverse
    selection}; the contract rate at a fixed offer rate identifies
    \textbf{moral hazard}.
  \item
    \textbf{Cole, Giné and Vickery (2017)} randomized rainfall
    insurance; insured farmers shifted into higher-return,
    rainfall-sensitive crops.
  \item
    \textbf{Duflo and Saez (2003)} paid \$20 to attend a benefits fair.
    Enrollment rose for treated employees \emph{and} untreated
    colleagues --- a clean randomization that violates SUTVA (Part 5).
  \end{itemize}
\item
  \textbf{What they share:} the researcher controls \textbf{an offer}
  --- a mailer, a price quote, a policy, an invitation. That is why
  field experiments cluster in \textbf{household finance, credit and
  insurance}.
\end{itemize}
\end{frame}

\begin{frame}{Why finance rarely gets to randomize}
\protect\phantomsection\label{why-finance-rarely-gets-to-randomize}
\begin{itemize}
\tightlist
\item
  Nobody will randomize leverage, M\&A, CEO pay, bankruptcy, or bank
  regulation:

  \begin{itemize}
  \tightlist
  \item
    stakes too high, units too big, ethics, general-equilibrium effects.
  \end{itemize}
\item
  So corporate finance and asset pricing live almost entirely in
  observational data.
\item
  \textbf{The rest of this course is a toolkit for \emph{earning}
  independence without a coin flip:}

  \begin{itemize}
  \tightlist
  \item
    Unconfoundedness / matching: independence \emph{conditional on
    observables}.
  \item
    RDD: independence \emph{near a threshold}.
  \item
    IV: independence \emph{of an instrument}.
  \item
    DiD / panel: independence \emph{of timing, in changes}.
  \end{itemize}
\item
  Every method = an argument that some comparison mimics the coin flip.
  Your job as a referee is to ask: \emph{whose coin, and who flipped
  it?}
\end{itemize}
\end{frame}

\section{5. SUTVA}\label{sutva}

\begin{frame}{The fine print under the potential outcomes}
\protect\phantomsection\label{the-fine-print-under-the-potential-outcomes}
\begin{itemize}
\tightlist
\item
  Writing \(Y_i^1, Y_i^0\) at all presumes the \textbf{Stable Unit
  Treatment Value Assumption}:
\end{itemize}

\begin{enumerate}
\tightlist
\item
  \textbf{No interference}: unit \(i\)'s potential outcomes do not
  depend on \emph{other units'} treatment status.
\item
  \textbf{No hidden versions of treatment}: ``treated'' means the same
  thing for every unit --- one treatment, one dose.
\end{enumerate}

\begin{itemize}
\tightlist
\item
  If SUTVA fails, \(Y_i^1\) is not even well-defined: bank \(i\)'s
  outcome under treatment depends on \emph{who else} was treated.
\end{itemize}
\end{frame}

\begin{frame}{Interference is the norm in finance}
\protect\phantomsection\label{interference-is-the-norm-in-finance}
\begin{itemize}
\tightlist
\item
  Units in finance \textbf{compete and interact}. Examples:

  \begin{itemize}
  \tightlist
  \item
    \textbf{Peer effects in capital structure} (Leary and Roberts 2014):
    firm \(i\)'s leverage responds to peers' leverage. Treating one firm
    treats its rivals.
  \item
    \textbf{Bailouts and competition}: a capital injection into bank
    \(i\) lets it bid for deposits and borrowers --- \emph{taking them
    from control banks}.
  \item
    \textbf{Banking DiDs}: treated banks' customers migrate to control
    banks; the control group's outcome moves \emph{because of} the
    treatment.
  \end{itemize}
\item
  If treatment hurts controls, the treated-control gap
  \textbf{overstates} the true effect; if it helps them (spillovers), it
  understates.
\end{itemize}
\end{frame}

\begin{frame}{``The control group is treated too''}
\protect\phantomsection\label{the-control-group-is-treated-too}
\begin{itemize}
\tightlist
\item
  In a competitive industry, there may be no untreated state of the
  world to observe:

  \begin{itemize}
  \tightlist
  \item
    Treatment effects estimated off treated-vs-control gaps are
    \textbf{relative} (partial-equilibrium) effects.
  \item
    The \textbf{aggregate} (general-equilibrium) effect can be smaller,
    zero, or of opposite sign.
  \end{itemize}
\item
  Example: a subsidy to treated firms raises their market share.
  Industry output may be unchanged --- pure reallocation. The DiD shows
  a large ``effect.''
\item
  Always ask: \emph{what happens to the market, not just the treated
  units?}
\item
  We will return to this in DiD (choosing clean controls) and in the
  reflection problem (peer effects).
\end{itemize}
\end{frame}

\begin{frame}{Hidden versions of treatment}
\protect\phantomsection\label{hidden-versions-of-treatment}
\begin{itemize}
\tightlist
\item
  SUTVA also requires the treatment to be \emph{one thing}.
\item
  Finance treatments are rarely uniform in dose:

  \begin{itemize}
  \tightlist
  \item
    ``Having a credit rating'' spans AAA to CCC.
  \item
    ``Being bailed out'' ranged from \$1bn to \$45bn in TARP.
  \item
    ``Adopting an anti-takeover provision'' bundles very different
    provisions.
  \end{itemize}
\item
  If doses differ and units select their dose, the single \(\delta_i\)
  hides a dose-response function --- and the estimand becomes an
  uninterpretable mixture.
\end{itemize}
\end{frame}

\begin{frame}{What to do about varying dose}
\protect\phantomsection\label{what-to-do-about-varying-dose}
\begin{itemize}
\tightlist
\item
  \textbf{Define} the treatment precisely. What exactly is the ``1'' in
  \(D_i = 1\)?
\item
  \textbf{Report} the \emph{distribution} of dose, not just a count of
  treated units.
\item
  If doses vary widely, do \textbf{not} collapse them into a binary.
  Instead, one of:

  \begin{itemize}
  \tightlist
  \item
    use the dose itself as the treatment and estimate a
    \textbf{dose-response function}, \(E[Y \mid \text{dose}]\);
  \item
    restrict the sample to a narrow dose band, so that ``treated'' means
    one thing;
  \item
    interact the treatment dummy with dose, and report the effect at
    several dose levels.
  \end{itemize}
\item
  Whichever you choose, state which estimand it delivers:
\end{itemize}

\begin{center}
\emph{``the effect of a \$1bn injection'' is not ``the effect of being bailed out.''}
\end{center}
\end{frame}

\section{6. Randomization Inference}\label{randomization-inference}

\begin{frame}{Why bother?}
\protect\phantomsection\label{why-bother}
\begin{itemize}
\tightlist
\item
  \textbf{Randomization-based inference}: the uncertainty in an estimate
  comes from the \textbf{random assignment of treatment}, not from
  hypothetical sampling out of a larger population.
\item
  The two approaches answer different questions:

  \begin{itemize}
  \tightlist
  \item
    A conventional standard error asks \emph{``what if I had drawn a
    different sample?''}
  \item
    Randomization inference asks \emph{``what if the treated had been
    differently?''}
  \end{itemize}
\item
  In an experiment, only the second thing actually happened. The sample
  is whoever was there; the \textbf{assignment} is what was random.
\item
  That shift buys you \emph{exact} \(p\)-values --- the probability that
  chance alone could have produced what you saw.
\item
  The idea is old: \textbf{R. A. Fisher}, early twentieth century.
\end{itemize}
\end{frame}

\begin{frame}{Three situations where it earns its keep}
\protect\phantomsection\label{three-situations-where-it-earns-its-keep}
\begin{enumerate}
\tightlist
\item
  \textbf{You have the population, not a sample.} With comprehensive or
  administrative data, a standard error answering ``what if I drew
  another sample?'' is not valuable. The real uncertainty is that we
  never see the counterfactual (Abadie, Diamond and Hainmueller 2010;
  Abadie et al.~2020).
\item
  \textbf{Few treated units.} Conventional inference leans on
  large-sample approximations. With a handful of treated units a single
  observation can swing the estimate, and conventional tests
  \textbf{over-reject} (Young, 2019).
\item
  \textbf{Transparency.} Placebo-based inference is easy to explain and
  hard to fudge.
\end{enumerate}

\vspace{0.4em}
\end{frame}

\begin{frame}{Two null hypotheses}
\protect\phantomsection\label{two-null-hypotheses}
\begin{itemize}
\tightlist
\item
  Standard (Neyman) null: the \textbf{average} effect is zero,
  \(E[\delta_i] = 0\).

  \begin{itemize}
  \tightlist
  \item
    Inference relies on asymptotics: large \(n\), estimated standard
    errors.
  \end{itemize}
\item
  Fisher's \textbf{sharp null}: the effect is zero \emph{for every
  unit}, \[
  H_0: \delta_i = Y_i^1 - Y_i^0 = 0 \quad \text{for all } i
  \]
\item
  The sharp null is \textbf{stronger}.
\item
  Its payoff is that it is \textbf{nonparametric} --- no Gaussian
  assumption, no large-sample approximation, no estimated standard
  error. It works at \(n = 8\).
\end{itemize}
\end{frame}

\begin{frame}{The sharp null fills in the missing half}
\protect\phantomsection\label{the-sharp-null-fills-in-the-missing-half}
\begin{itemize}
\tightlist
\item
  The fundamental problem of causal inference is that we never see both
  potential outcomes. Under the sharp null, \textbf{we do}.
\item
  If \(\delta_i = Y_i^1 - Y_i^0 = 0\), then \(Y_i^1 = Y_i^0\), and both
  equal the value we observed.
\end{itemize}

\begin{center}
\small
\begin{tabular}{c|cc|cc||cc}
\toprule
 & \multicolumn{4}{c||}{What we observe} & \multicolumn{2}{c}{Under $H_0: \delta_i = 0$} \\
Bank & $D$ & $Y$ & $Y^0$ & $Y^1$ & $Y^0$ & $Y^1$ \\
\midrule
1 & 0 & 2 & 2 & ? & 2 & 2 \\
2 & 1 & 6 & ? & 6 & 6 & 6 \\
3 & 1 & 7 & ? & 7 & 7 & 7 \\
$\vdots$ & & & & & & \\
10 & 0 & 8 & 8 & ? & 8 & 8 \\
\bottomrule
\end{tabular}
\end{center}
\end{frame}

\begin{frame}{It is not the bootstrap}
\protect\phantomsection\label{it-is-not-the-bootstrap}
\begin{itemize}
\tightlist
\item
  The two get confused because both produce a distribution by
  resampling.
\item
  \textbf{Bootstrap.} Resample \emph{observations}, with replacement,
  from your sample. The resulting uncertainty is about \textbf{which
  units ended up in your sample} --- sampling uncertainty.
\item
  \textbf{Randomization inference.} Hold the units fixed and reshuffle
  \textbf{who was treated}. The uncertainty is about \textbf{which units
  got the treatment} --- assignment uncertainty.
\item
  Different questions, different null, different interpretation.
\end{itemize}
\end{frame}

\begin{frame}{Six steps}
\protect\phantomsection\label{six-steps}
Following Cunningham (\emph{The Remix}, §4.7):

\begin{enumerate}
\tightlist
\item
  Choose the \textbf{sharp null}.
\item
  Construct a \textbf{test statistic} \(t(D, Y)\).
\item
  Pick a \textbf{different} treatment assignment \(\tilde D\).
\item
  Compute the test statistic for that pretend assignment.
\item
  \textbf{Cycle} over step 3 for all possible assignments.
\item
  Divide the \textbf{rank} of the true statistic by the number of trials
  --- the exact \(p\)-value.
\end{enumerate}
\end{frame}

\begin{frame}{Step 2: what makes a good test statistic}
\protect\phantomsection\label{step-2-what-makes-a-good-test-statistic}
\begin{itemize}
\tightlist
\item
  A test statistic \(t(D, Y)\) is just a \textbf{scalar} built from the
  assignment and the outcomes.
\item
  The requirement: it should take \textbf{extreme values when the sharp
  null is false}, and extreme values should be \textbf{rare when it is
  true}.
\item
  The obvious choice is the one we already know --- the simple
  difference in outcomes, \[
  t(D, Y) = \bar{Y}_{\text{treated}} - \bar{Y}_{\text{control}} = SDO
  \]
\item
  For the ten banks, the lottery gave injections to banks
  \(2, 3, 5, 8, 9\): \[
  SDO = \frac{6+7+6+7+7}{5} - \frac{2+2+6+7+8}{5} = 6.6 - 5.0 = 1.6
  \]
\end{itemize}
\end{frame}

\begin{frame}{Steps 3--5: the permutation loop}
\protect\phantomsection\label{steps-35-the-permutation-loop}
\begin{itemize}
\tightlist
\item
  Under the sharp null the outcomes \(Y\) \textbf{do not change} ---
  only the labels move.
\item
  Suppose banks \(1, 3, 5, 8, 9\) had been drawn instead: \[
  \tilde t = \frac{2+7+6+7+7}{5} - \frac{6+2+6+7+8}{5} = 5.8 - 5.8 = 0.0
  \]
\item
  Or banks \(1, 2, 4, 6, 7\): \[
  \tilde t = \frac{2+6+2+6+7}{5} - \frac{7+6+7+7+8}{5} = 4.6 - 7.0 = -2.4
  \]
\item
  Store each one and move on. With five treated out of ten there are
  \(\binom{10}{5} = 252\) such assignments.
\item
  Each \(\tilde t\) is a number the world \emph{could have produced} if
  the treatment did nothing. Together they are the \textbf{reference
  distribution}.
\end{itemize}
\end{frame}

\begin{frame}{Step 6: the exact \(p\)-value}
\protect\phantomsection\label{step-6-the-exact-p-value}
\[
\Pr\big(|t(\tilde D, Y)| \geq |t(D, Y)| \;\big|\; \delta_i = 0 \;\, \forall i\big) = \frac{\sum_{\tilde D \in \Omega} I\big(|t(\tilde D, Y)| \geq |t(D, Y)|\big)}{K}
\]

\begin{itemize}
\tightlist
\item
  \textbf{what share of the pretend worlds look at least as extreme as
  the real one?}
\item
  \(\Omega\) is the set of assignments the randomization could have
  produced and \(K = |\Omega|\). Here \(K = 252\).
\item
  ``\(\delta = 0\)'' would name the Neyman null (\emph{The Remix})
\item
  \textbf{``Exact''} is meant literally. This is not an approximation to
  a limiting distribution --- it is a finite count.
\end{itemize}
\end{frame}

\begin{frame}[fragile]{Randomization inference on ten banks}
\protect\phantomsection\label{randomization-inference-on-ten-banks}
\begin{Shaded}
\begin{Highlighting}[]
\NormalTok{y }\OtherTok{\textless{}{-}} \FunctionTok{c}\NormalTok{(}\DecValTok{2}\NormalTok{, }\DecValTok{6}\NormalTok{, }\DecValTok{7}\NormalTok{, }\DecValTok{2}\NormalTok{, }\DecValTok{6}\NormalTok{, }\DecValTok{6}\NormalTok{, }\DecValTok{7}\NormalTok{, }\DecValTok{7}\NormalTok{, }\DecValTok{7}\NormalTok{, }\DecValTok{8}\NormalTok{)     }\CommentTok{\# observed lending growth}
\NormalTok{treated }\OtherTok{\textless{}{-}} \FunctionTok{c}\NormalTok{(}\DecValTok{2}\NormalTok{, }\DecValTok{3}\NormalTok{, }\DecValTok{5}\NormalTok{, }\DecValTok{8}\NormalTok{, }\DecValTok{9}\NormalTok{)              }\CommentTok{\# banks drawn by lot}
\NormalTok{sdo\_obs }\OtherTok{\textless{}{-}} \FunctionTok{mean}\NormalTok{(y[treated]) }\SpecialCharTok{{-}} \FunctionTok{mean}\NormalTok{(y[}\SpecialCharTok{{-}}\NormalTok{treated])}

\NormalTok{perms }\OtherTok{\textless{}{-}} \FunctionTok{combn}\NormalTok{(}\DecValTok{10}\NormalTok{, }\DecValTok{5}\NormalTok{)                     }\CommentTok{\# all 252 possible assignments}
\NormalTok{sdo\_all }\OtherTok{\textless{}{-}} \FunctionTok{apply}\NormalTok{(perms, }\DecValTok{2}\NormalTok{, }\ControlFlowTok{function}\NormalTok{(tr) }\FunctionTok{mean}\NormalTok{(y[tr]) }\SpecialCharTok{{-}} \FunctionTok{mean}\NormalTok{(y[}\SpecialCharTok{{-}}\NormalTok{tr]))}
\NormalTok{p\_exact }\OtherTok{\textless{}{-}} \FunctionTok{mean}\NormalTok{(}\FunctionTok{abs}\NormalTok{(sdo\_all) }\SpecialCharTok{\textgreater{}=} \FunctionTok{abs}\NormalTok{(sdo\_obs))}
\FunctionTok{round}\NormalTok{(}\FunctionTok{c}\NormalTok{(}\AttributeTok{SDO =}\NormalTok{ sdo\_obs, }\AttributeTok{exact\_p =}\NormalTok{ p\_exact), }\DecValTok{3}\NormalTok{)}
\end{Highlighting}
\end{Shaded}

\begin{verbatim}
    SDO exact_p 
  1.600   0.397 
\end{verbatim}

\begin{itemize}
\tightlist
\item
  \(p = 100/252 = 0.397\). With \(n = 10\), a difference of 1.6pp is
  \textbf{not} distinguishable from luck of the draw --- and we know
  this \emph{exactly}, not asymptotically.
\end{itemize}
\end{frame}

\begin{frame}[fragile]{The assumption hiding in \texttt{combn(10,\ 5)}}
\protect\phantomsection\label{the-assumption-hiding-in-combn10-5}
\begin{itemize}
\tightlist
\item
  The reference set must match the \textbf{actual randomization design}.
\item
  \texttt{combn(10,\ 5)} presumes a completely randomized experiment
  with a \textbf{fixed} five treated --- exactly five drawn by lot.
\item
  Had each bank instead been \textbf{independently coin-flipped}, the
  correct reference set would be all \(2^{10} = 1024\) assignments, and
  the \(p\)-value would differ.
\item
  Get this wrong and the \(p\)-value is not ``approximately right.'' It
  answers a different question.
\end{itemize}

\begin{center}
\emph{The permutation scheme must mirror the physical randomization.}
\end{center}
\end{frame}

\begin{frame}{Revisiting step 2: the menu of test statistics}
\protect\phantomsection\label{revisiting-step-2-the-menu-of-test-statistics}
\begin{itemize}
\tightlist
\item
  \(SDO\) is the default, and a good one \textbf{when effects are
  additive and outliers few}.
\item
  \textbf{Log difference} --- for skewed outcomes and multiplicative
  effects: \[
  T_{\log} = \frac{1}{N_T}\sum_i D_i \ln Y_i - \frac{1}{N_C}\sum_i (1 - D_i)\ln Y_i
  \]
\item
  \textbf{Quantiles} --- robust to outliers in either tail: \[
  T_{\text{median}} = \big|\,\text{median}(Y_T) - \text{median}(Y_C)\,\big|
  \]
\item
  \textbf{Normalised ranks} --- discard magnitudes, keep the ordering.
  Good with small samples or heavy tails.
\item
  \textbf{Kolmogorov--Smirnov} --- looks at the whole distribution, not
  a single summary.
\item
  \textbf{None of these changes the null or the validity of the test}.
  The choice changes only \textbf{power}. So choose \textbf{before} you
  look.
\end{itemize}
\end{frame}

\begin{frame}{Revisiting step 5: when ``all assignments'' is impossible}
\protect\phantomsection\label{revisiting-step-5-when-all-assignments-is-impossible}
\begin{itemize}
\tightlist
\item
  Ten banks give \(252\) assignments. In \emph{The Remix}, Thornton's
  (2008) HIV-testing experiment had \(2{,}901\) participants,
  \(2{,}222\) of them incentivised. The number of assignments is \[
  \binom{2901}{2222} \approx 6.15 \times 10^{680}
  \]
\item
  So we \textbf{sample} the reference distribution: draw a random
  assignment, recompute the statistic, repeat \(1{,}000\) to
  \(10{,}000\) times.
\item
  The \(p\)-value is then \textbf{approximate} rather than exact --- but
  the error is sampling error you control with more draws, not an
  asymptotic assumption you must defend.
\item
  This is what you will actually do. Any dataset of a reasonable size
  makes full enumeration hopeless.
\end{itemize}
\end{frame}

\begin{frame}{Where you will see this again}
\protect\phantomsection\label{where-you-will-see-this-again}
\begin{itemize}
\tightlist
\item
  Randomization inference is not a curiosity:

  \begin{itemize}
  \tightlist
  \item
    \textbf{Synthetic control}: placebo tests permute the treatment
    across donor units --- that \emph{is} randomization inference.
  \item
    \textbf{Small-N settings}: one treated state, one regulatory event,
    a handful of treated firms --- common in finance, hopeless for
    asymptotics.
  \item
    Cluster-robust inference with few clusters: permutation tests as the
    honest fallback.
  \end{itemize}
\item
  The deeper point: \emph{inference should mirror the assignment
  mechanism.} Where the mechanism is uncertain, so is the \(p\)-value.
\end{itemize}
\end{frame}

\section{Wrapping Up}\label{wrapping-up}

\begin{frame}{Summary {[}Part 1{]}}
\protect\phantomsection\label{summary-part-1}
\begin{itemize}
\tightlist
\item
  Causality is defined by \textbf{counterfactuals}; the fundamental
  problem is that one potential outcome is always missing.
\item
  Estimands differ: \textbf{ATE, ATT, ATU} answer different policy
  questions. Know which one your design identifies.
\item
  The naive comparison decomposes as \[
  SDO = ATE + \text{selection bias} + (1-\pi)(ATT - ATU)
  \]
\item
  Optimizing agents generate selection \textbf{on gains} --- the Perfect
  Regulator gets the sign wrong \emph{because} it allocates well.
\item
  More data does not fix a biased assignment mechanism.
\end{itemize}
\end{frame}

\begin{frame}{Summary {[}Part 2{]}}
\protect\phantomsection\label{summary-part-2}
\begin{itemize}
\tightlist
\item
  \textbf{Randomization} kills both bias terms by making
  \((Y^1, Y^0) \indep D\). Balance tables are its testable symptom.
\item
  \textbf{SUTVA} fails easily in finance: competition, peer effects, and
  heterogeneous doses. Ask what happens to the market, not just the
  treated.
\item
  \textbf{Randomization inference} turns the assignment mechanism into a
  test. Under \textbf{Fisher's sharp null} every counterfactual is
  filled in, so the \(p\)-value is an exact count over the assignments
  the lottery could have produced --- no asymptotics, at any sample
  size.
\item
  The \textbf{test statistic is a choice}: means, logs, quantiles,
  ranks, or the whole distribution (Kolmogorov--Smirnov). Make it before
  you look, and remember that a treatment can move the \emph{spread}
  without moving the mean.
\item
  Next lecture: DAGs --- a graphical language for deciding \emph{what to
  control for, and what not to}.
\end{itemize}
\end{frame}

\begin{frame}{References}
\protect\phantomsection\label{references}
\footnotesize

\begin{itemize}
\tightlist
\item
  Cunningham, S. \emph{Causal Inference: The Remix}, Ch. 4.
\item
  Holland, P. (1986). Statistics and Causal Inference. \emph{JASA} 81,
  945--960.
\item
  Abadie, A., A. Diamond and J. Hainmueller (2010). Synthetic Control
  Methods for Comparative Case Studies. \emph{JASA} 105, 493--505.
\item
  Abadie, A., S. Athey, G. Imbens and J. Wooldridge (2020).
  Sampling-Based versus Design-Based Uncertainty in Regression Analysis.
  \emph{Econometrica} 88, 265--296.
\item
  Fisher, R. A. (1935). \emph{The Design of Experiments}. Oliver and
  Boyd.
\item
  Thornton, R. (2008). The Demand for, and Impact of, Learning HIV
  Status. \emph{AER} 98, 1829--1863.
\item
  Young, A. (2019). Channeling Fisher: Randomization Tests and the
  Statistical Insignificance of Seemingly Significant Experimental
  Results. \emph{QJE} 134, 557--598.
\item
  Blake, T., C. Nosko, and S. Tadelis (2015). Consumer Heterogeneity and
  Paid Search Effectiveness: A Large-Scale Field Experiment.
  \emph{Econometrica} 83, 155--174.
\item
  Karlan, D. and J. Zinman (2009). Observing Unobservables: Identifying
  Information Asymmetries with a Consumer Credit Field Experiment.
  \emph{Econometrica} 77, 1993--2008.
\item
  Cole, S., X. Giné, and J. Vickery (2017). How Does Risk Management
  Influence Production Decisions? \emph{Review of Financial Studies} 30,
  1935--1970.
\item
  Duflo, E. and E. Saez (2003). The Role of Information and Social
  Interactions in Retirement Plan Decisions. \emph{Quarterly Journal of
  Economics} 118, 815--842.
\item
  Leary, M. and M. Roberts (2014). Do Peer Firms Affect Corporate
  Financial Policy? \emph{Journal of Finance} 69, 139--178.
\end{itemize}
\end{frame}




\end{document}
